% 数学ノート集・共通テンプレート v2（積分論2026年度第1部の版面を参考）
% A5判。章内で定義・定理・例・注意等の番号を共通化。
% コンパイル： platex -interaction=nonstopmode math_note_template_v2.tex (2回)
%              dvipdfmx math_note_template_v2.dvi
\documentclass[a5paper,10pt,onecolumn]{jsarticle}
\usepackage[dvipdfmx,top=22mm,bottom=20mm,inner=19mm,outer=18mm,headheight=13pt,headsep=6mm]{geometry}
\usepackage{amsmath,amssymb,amsthm,mathtools}
\usepackage{enumitem,fancyhdr,titlesec}
\usepackage[dvipdfmx,bookmarks=true,colorlinks=false]{hyperref}
\usepackage{pxjahyper}

% ----- 差し替える文書情報 -----
\newcommand{\NoteTitle}{数値最適化法と収束の定理}
\newcommand{\NoteSeries}{数学ノート集}

% ----- 文字・組版 -----
% 参考PDFの質素な明朝＋Computer Modernの印象を引き継ぎながら，
% A5判で窮屈にならないよう本文の行間を少し確保する．
\linespread{1.13}
\setlength{\parindent}{1zw}
\setlength{\parskip}{0pt}
\setlength{\abovedisplayskip}{6pt plus 2pt minus 2pt}
\setlength{\belowdisplayskip}{6pt plus 2pt minus 2pt}
\setlength{\abovedisplayshortskip}{3pt plus 2pt minus 1pt}
\setlength{\belowdisplayshortskip}{4pt plus 2pt minus 1pt}
\setlist[enumerate]{topsep=3pt,itemsep=1pt,parsep=0pt,leftmargin=2.4zw}
\titleformat{\section}{\normalfont\normalsize\bfseries}{\thesection}{0.8zw}{}
\titlespacing*{\section}{0pt}{2.0ex plus .4ex}{.9ex}
\titleformat{\subsection}{\normalfont\normalsize\bfseries}{\thesubsection}{0.7zw}{}
\titlespacing*{\subsection}{0pt}{1.5ex plus .3ex}{.6ex}
\raggedbottom

% ----- 柱とページ番号：ページ上部、罫線なし -----
\fancypagestyle{mathnote}{%
  \fancyhf{}%
  \fancyhead[L]{\footnotesize \NoteSeries\ （\NoteTitle）}%
  \fancyhead[R]{\footnotesize \thepage}%
  \renewcommand{\headrulewidth}{0pt}%
  \renewcommand{\footrulewidth}{0pt}%
}
\pagestyle{mathnote}

% ----- 節内通し番号の定義・定理群 -----
\newtheoremstyle{mathnotetheorem}
 {4.5pt}{4.5pt}{\normalfont}{}%
 {\bfseries}{}{0.65zw}%
 {\thmname{#1}\thmnumber{ #2}\thmnote{（#3）}}
\theoremstyle{mathnotetheorem}
\newtheorem{theorem}{定理}[section]
\newtheorem{definition}[theorem]{定義}
\newtheorem{lemma}[theorem]{補題}
\newtheorem{proposition}[theorem]{命題}
\newtheorem{corollary}[theorem]{系}
\newtheorem{example}[theorem]{例}
\newtheorem{remark}[theorem]{注意}
\renewcommand{\proofname}{証明}
\renewcommand{\qedsymbol}{$\square$}
\makeatletter
\renewenvironment{proof}[1][\proofname]{%
  \par\pushQED{\qed}\normalfont\topsep4\p@\@plus2\p@\relax
  \trivlist\item[\hskip\labelsep\bfseries #1]\ignorespaces
}{\popQED\endtrivlist\@endpefalse}
\makeatother

% ----- 使うノートで適宜拡張する記号 -----
\newcommand{\R}{\mathbb{R}}
\newcommand{\norm}[1]{\lVert #1\rVert}
\newcommand{\abs}[1]{\lvert #1\rvert}

\begin{document}
\section*{読了後にできること}
本稿を読むための前提は，勾配降下法，収束とスペクトル半径．
本稿の具体的な到達課題は，簡単な二次関数で刻み幅の収束条件を厳密に求める．


\section{滑らかな凸関数と降下補題}

本稿では $\mathbb R^n$ にEuclid内積とノルムを入れる．
微分可能な凸関数の一次不等式
\[
 f(y)\ge f(x)+\langle\nabla f(x),y-x\rangle
\]
を既知とする．

\begin{definition}
微分可能な関数 $f:\mathbb R^n\to\mathbb R$ の
勾配が\emph{$L$-Lipschitz連続}であるとは，
$L>0$ が存在して
\[
 \norm{\nabla f(x)-\nabla f(y)}
 \le L\norm{x-y}\qquad(x,y\in\mathbb R^n)
\]
を満たすことをいう．
\end{definition}

\begin{lemma}[降下補題]\label{lem:descent}
$f$ の勾配が $L$-Lipschitz連続なら
\[
 f(y)\le f(x)+\langle\nabla f(x),y-x\rangle
       +\frac L2\norm{y-x}^2
\]
が任意の $x,y$ に対して成立する．
\end{lemma}
\begin{proof}
$h=y-x$ と置く．
微積分の基本定理から
\[
 f(y)-f(x)
 =\int_0^1\langle\nabla f(x+th),h\rangle\,dt.
\]
従って
\begin{align*}
 &f(y)-f(x)-\langle\nabla f(x),h\rangle\\
 &\quad=\int_0^1
 \langle\nabla f(x+th)-\nabla f(x),h\rangle\,dt\\
 &\quad\le\int_0^1Lt\norm h^2\,dt
 =\frac L2\norm h^2.
\end{align*}
\end{proof}

\begin{definition}
微分可能な関数 $f$ が $\mu$-\emph{強凸}
（$\mu>0$）であるとは，
\[
 f(y)\ge f(x)+\langle\nabla f(x),y-x\rangle
 +\frac\mu2\norm{y-x}^2
\]
が任意の $x,y$ に対して成立することをいう．
\end{definition}

\begin{proposition}[強凸関数の最小点]\label{prop:strong-min}
$f$ が微分可能で $\mu$-強凸なら，
$f$ はただ一つの最小点 $x_*$ を持ち，
$\nabla f(x_*)=0$ を満たす．
\end{proposition}
\begin{proof}
強凸性を $x=0$ に適用すると
\[
 f(y)\ge f(0)+\langle\nabla f(0),y\rangle
 +\frac\mu2\norm y^2
\]
となる．右辺は $\norm y\to\infty$ で
$\infty$ へ向かうので $f$ は強制的である．
連続性とHeine--Borelの定理によって最小点が存在する．
微分可能な関数の内点の最小点では $\nabla f=0$ である．
二つの最小点 $x_*,y_*$ があれば強凸性から
\[
 f(y_*)\ge f(x_*)+\frac\mu2\norm{y_*-x_*}^2
\]
となる．両辺の関数値は等しいので $y_*=x_*$ である．
\end{proof}

\section{勾配降下法}

\begin{definition}
勾配が $L$-Lipschitz連続な $f$ に対して，
\[
 x_{k+1}=x_k-\frac1L\nabla f(x_k)
\]
で定める列を固定歩長の\emph{勾配降下法}という．
\end{definition}

\begin{lemma}[一歩の降下]\label{lem:one-step}
勾配降下列について
\[
 f(x_{k+1})\le f(x_k)
 -\frac1{2L}\norm{\nabla f(x_k)}^2
\]
が成り立つ．
\end{lemma}
\begin{proof}
降下補題において
$y=x_k-\frac1L\nabla f(x_k)$ と置く．
\begin{align*}
 f(x_{k+1})
 &\le f(x_k)
 -\frac1L\norm{\nabla f(x_k)}^2
 +\frac L2\frac1{L^2}\norm{\nabla f(x_k)}^2\\
 &=f(x_k)-\frac1{2L}\norm{\nabla f(x_k)}^2.
\end{align*}
\end{proof}

\begin{theorem}[凸関数に対する収束率]\label{thm:convex-rate}
$f:\mathbb R^n\to\mathbb R$ を微分可能な凸関数で，
勾配は $L$-Lipschitz連続とする．
最小点 $x_*$ が存在すると仮定する．
勾配降下列に対して任意の $N\ge1$ で
\[
 0\le f(x_N)-f(x_*)
 \le\frac{L\norm{x_0-x_*}^2}{2N}
\]
が成り立つ．
\end{theorem}
\begin{proof}
凸性の一次不等式から
\[
 f(x_k)-f(x_*)
 \le\langle\nabla f(x_k),x_k-x_*\rangle.
\]
補題\ref{lem:one-step}と合わせると
\[
 f(x_{k+1})-f(x_*)
 \le\langle\nabla f(x_k),x_k-x_*\rangle
 -\frac1{2L}\norm{\nabla f(x_k)}^2.
\]
一方，更新式から
\begin{align*}
 \norm{x_{k+1}-x_*}^2
 &=\norm{x_k-x_*}^2
 -\frac2L\langle\nabla f(x_k),x_k-x_*\rangle\\
 &\qquad+\frac1{L^2}\norm{\nabla f(x_k)}^2.
\end{align*}
両式を比較すると
\[
 f(x_{k+1})-f(x_*)
 \le\frac L2
 \bigl(\norm{x_k-x_*}^2-\norm{x_{k+1}-x_*}^2\bigr).
\]
$k=0,\dots,N-1$ について足すと
\[
 \sum_{k=1}^N(f(x_k)-f(x_*))
 \le\frac L2\norm{x_0-x_*}^2.
\]
補題\ref{lem:one-step}から $f(x_k)$ は単調非増加なので
左辺は $N(f(x_N)-f(x_*))$ 以上である．
従って結論が従う．
\end{proof}

\begin{corollary}[勾配の消失]\label{cor:grad-zero}
$f$ が下に有界で勾配は $L$-Lipschitz連続なら，
勾配降下列は
\[
 \sum_{k=0}^\infty\norm{\nabla f(x_k)}^2<\infty,
 \qquad \norm{\nabla f(x_k)}\to0
\]
を満たす．
\end{corollary}
\begin{proof}
補題\ref{lem:one-step}を $k=0,\dots,N-1$ について足すと
\[
 \frac1{2L}\sum_{k=0}^{N-1}\norm{\nabla f(x_k)}^2
 \le f(x_0)-f(x_N)
 \le f(x_0)-\inf f.
\]
$N\to\infty$ とすれば非負項級数が収束する．
従って一般項は零に収束する．
\end{proof}

\section{強凸性による線形収束}

\begin{lemma}[強凸関数の勾配下界]\label{lem:pl}
$f$ を微分可能な $\mu$-強凸関数，
$x_*$ をその最小点とする．
このとき
\[
 \norm{\nabla f(x)}^2
 \ge2\mu\bigl(f(x)-f(x_*)\bigr)
\]
が成り立つ．
\end{lemma}
\begin{proof}
強凸性により任意の $y$ に対して
\[
 f(y)\ge f(x)+\langle\nabla f(x),y-x\rangle
 +\frac\mu2\norm{y-x}^2.
\]
右辺は $y=x-\mu^{-1}\nabla f(x)$ のとき最小であり，
その最小値は
$f(x)-\norm{\nabla f(x)}^2/(2\mu)$ となる．
特に $y=x_*$ として
\[
 f(x_*)\ge f(x)-\frac1{2\mu}\norm{\nabla f(x)}^2
\]
なので整理して結論を得る．
\end{proof}

\begin{theorem}[強凸の場合の線形収束]\label{thm:linear-rate}
$f$ を微分可能な $\mu$-強凸関数とし，
その勾配は $L$-Lipschitz連続とする．
このとき $0<\mu\le L$ であり，
勾配降下法について
\[
 f(x_k)-f(x_*)
 \le\left(1-\frac\mu L\right)^k
 \bigl(f(x_0)-f(x_*)\bigr)
\]
が成り立つ．
\end{theorem}
\begin{proof}
強凸性と降下補題を同じ二点について比較すると，
その差から $\mu\le L$ が従う
（$x\ne y$ として一次項を打ち消せばよい）．
補題\ref{lem:one-step}と補題\ref{lem:pl}から
\begin{align*}
 f(x_{k+1})-f(x_*)
 &\le f(x_k)-f(x_*)-\frac1{2L}\norm{\nabla f(x_k)}^2\\
 &\le\left(1-\frac\mu L\right)
       (f(x_k)-f(x_*)).
\end{align*}
これを帰納的に繰り返して結論を得る．
\end{proof}

\section{Newton法と局所二次収束}

\begin{definition}
$f:\mathbb R^n\to\mathbb R$ が二回連続微分可能で，
$\nabla^2f(x_k)$ が可逆であるとき，
\[
 x_{k+1}
 =x_k-\nabla^2f(x_k)^{-1}\nabla f(x_k)
\]
を\emph{Newton法}という．
\end{definition}

\begin{theorem}[Newton法の局所二次収束]\label{thm:newton}
$f\in C^2(\mathbb R^n)$ とし，
$\nabla f(x_*)=0$ かつ
$H_*=\nabla^2f(x_*)$ は正定値とする．
ある $r_0>0$ と $M>0$ が存在して，
$B(x_*,r_0)$ 上で
\[
 \norm{\nabla^2f(x)-\nabla^2f(y)}
 \le M\norm{x-y}
\]
が成り立つと仮定する．
このとき $x_*$ の十分小さい近傍内から開始するNewton列は
well-definedで $x_*$ に収束し，
ある定数 $C>0$ について
\[
 \norm{x_{k+1}-x_*}\le C\norm{x_k-x_*}^2
\]
を満たす．
\end{theorem}
\begin{proof}
$H_*$ は可逆であり，
Hessianは連続なので，
十分小さい閉球 $\overline{B(x_*,r)}$
において全ての $\nabla^2f(x)$ は可逆で，
その逆行列のノルムは一様に $B>0$ で抑えられる．
実際
\[
 \nabla^2f(x)=H_*
 \bigl(I+H_*^{-1}(\nabla^2f(x)-H_*)\bigr)
\]
と書き，括弧内にNeumann級数を適用すればよい．
$e=x-x_*$ と置く．
$\nabla f(x_*)=0$ と微積分の基本定理から
\[
 \nabla f(x)=\int_0^1\nabla^2f(x_*+te)e\,dt.
\]
Newton更新点を $N(x)$ と書くと
\begin{align*}
 N(x)-x_*
 &=e-\nabla^2f(x)^{-1}\nabla f(x)\\
 &=\nabla^2f(x)^{-1}
 \int_0^1\bigl(\nabla^2f(x)
 -\nabla^2f(x_*+te)\bigr)e\,dt.
\end{align*}
HessianのLipschitz条件から
\[
 \norm{N(x)-x_*}
 \le BM\norm e^2\int_0^1(1-t)\,dt
 =\frac{BM}{2}\norm e^2.
\]
$C=BM/2$ と置き，
必要なら $r$ をさらに小さくして $Cr<1$ とする．
$\norm{x_0-x_*}<r$ なら上の評価から
$\norm{N(x_0)-x_*}<r$ であり，
帰納的に全ての反復点は球内に留まる．
さらに
$\norm{x_{k+1}-x_*}\le Cr\norm{x_k-x_*}$
なので誤差は零へ収束する．
同時に二次誤差評価も成立する．
\end{proof}

\section{射影勾配法}

\begin{theorem}[閉凸集合への射影]\label{thm:projection}
$C\subset\mathbb R^n$ を空でない閉凸集合とする．
任意の $x\in\mathbb R^n$ に対して
\[
 P_C(x)=\operatorname*{argmin}_{z\in C}\norm{x-z}
\]
は一意に存在する．
また
\[
 \langle x-P_C(x),z-P_C(x)\rangle\le0
 \qquad(z\in C)
\]
が成り立つ．
\end{theorem}
\begin{proof}
$z_0\in C$ を固定する．
距離を最小化する点は
$C\cap\overline{B(x,\norm{x-z_0})}$ の中に存在すれば十分である．
この集合は非空・閉・有界なのでコンパクトであり，
連続関数 $z\mapsto\norm{x-z}$ は最小値を達成する．
二つの異なる最小点 $u,v$ があるとすると，
中点 $(u+v)/2\in C$ に対して
\[
 \norm{x-\frac{u+v}2}^2
 =\frac12\norm{x-u}^2+\frac12\norm{x-v}^2
  -\frac14\norm{u-v}^2
\]
となり，最小性に矛盾する．
従って射影は一意である．
$u=P_C(x)$ と置くと，
任意の $z\in C$ と $0<t<1$ に対して
$u+t(z-u)\in C$ なので
$\norm{x-u}^2\le\norm{x-u-t(z-u)}^2$ である．
展開して $t$ で割り $t\downarrow0$ とすると
$\langle x-u,z-u\rangle\le0$ を得る．
\end{proof}

\begin{proposition}[射影の非拡大性]\label{prop:projection-lipschitz}
$P_C$ は $1$-Lipschitz連続である．
\end{proposition}
\begin{proof}
$u=P_C(x)$，$v=P_C(y)$ と置く．
定理\ref{thm:projection}から
$\langle x-u,v-u\rangle\le0$，
$\langle y-v,u-v\rangle\le0$ である．
両辺を整理して足すと
\[
 \norm{u-v}^2\le\langle x-y,u-v\rangle
 \le\norm{x-y}\norm{u-v}.
\]
$u=v$ なら明らかであり，そうでなければ
$\norm{u-v}$ で割って
$\norm{P_C(x)-P_C(y)}\le\norm{x-y}$ を得る．
\end{proof}

\begin{definition}
$C\subset\mathbb R^n$ を空でない閉凸集合，
$f$ を微分可能で勾配が $L$-Lipschitz連続な凸関数とする．
反復
\[
 x_{k+1}
 =P_C\left(x_k-\frac1L\nabla f(x_k)\right)
\]
を\emph{射影勾配法}という．
\end{definition}

\begin{theorem}[射影勾配法の収束率]\label{thm:proj-rate}
上の仮定のもとで $f$ は $C$ 上で最小点 $x_*$ を持つとする．
$x_0\in C$ として射影勾配法を実行すると，
\[
 f(x_N)-f(x_*)\le
 \frac{L\norm{x_0-x_*}^2}{2N}\qquad(N\ge1)
\]
が成り立つ．
\end{theorem}
\begin{proof}
$z=x_k$，$u=x_{k+1}$ と置く．
射影の特徴づけを
$x=z-\nabla f(z)/L$ に適用すると，
任意の $y\in C$ に対して
\[
 \langle\nabla f(z)+L(u-z),y-u\rangle\ge0
\]
となる．
これを整理し，恒等式
\[
 2\langle u-z,y-u\rangle
 =\norm{y-z}^2-\norm{y-u}^2-\norm{u-z}^2
\]
を用いると
\begin{align*}
 \langle\nabla f(z),u-y\rangle
 &\le\frac L2\bigl(
 \norm{y-z}^2-\norm{y-u}^2-\norm{u-z}^2
 \bigr).
\end{align*}
一方，降下補題と凸性から
\[
 f(u)-f(y)
 \le\langle\nabla f(z),u-y\rangle
 +\frac L2\norm{u-z}^2.
\]
両式を足すと
\[
 f(u)-f(y)
 \le\frac L2(\norm{y-z}^2-\norm{y-u}^2).
 \tag{*}
\]
$y=x_k=z$ とすると $f(x_{k+1})\le f(x_k)$ だから
関数値は単調非増加である．
$y=x_*$ として (*) を $k=0,\dots,N-1$ について足すと
\[
 \sum_{k=1}^N(f(x_k)-f(x_*))
 \le\frac L2\norm{x_0-x_*}^2.
\]
単調性により左辺は
$N(f(x_N)-f(x_*))$ 以上であるから，
結論を得る．
\end{proof}

\section{近接作用素と反例}

\begin{definition}
$f:\mathbb R^n\to\mathbb R$ を連続な凸関数，
$\tau>0$ とする．
\[
 \operatorname{prox}_{\tau f}(x)
 =\operatorname*{argmin}_{z\in\mathbb R^n}
 \left\{f(z)+\frac1{2\tau}\norm{z-x}^2\right\}
\]
を\emph{近接作用素}という．
\end{definition}

\begin{proposition}[近接作用素の存在と非拡大性]
\label{prop:prox}
$\operatorname{prox}_{\tau f}(x)$ は
全ての $x$ に対して一意に存在する．
また
\[
 \norm{\operatorname{prox}_{\tau f}(x)
       -\operatorname{prox}_{\tau f}(y)}
 \le\norm{x-y}
\]
が成り立つ．
\end{proposition}
\begin{proof}
凸関数の一次下界
$f(z)\ge f(0)+\langle v,z\rangle$ を取れる．
連続な有限値凸関数は劣勾配を持つので，
例えば $v\in\partial f(0)$ を取ればよい．
従って目的関数は二次項によって強制的であり，
連続なので最小点が存在する．
二次項が厳密凸で $f$ が凸だから，最小点は一意である．

$u=\operatorname{prox}_{\tau f}(x)$，
$v=\operatorname{prox}_{\tau f}(y)$ と置く．
最小性の一次条件から
\[
 \frac{x-u}{\tau}\in\partial f(u),\qquad
 \frac{y-v}{\tau}\in\partial f(v)
\]
を得る．
劣微分の単調性を示すため，
$p\in\partial f(u)$，$q\in\partial f(v)$ に対し
\[
 f(v)\ge f(u)+\langle p,v-u\rangle,\qquad
 f(u)\ge f(v)+\langle q,u-v\rangle
\]
を足すと
$\langle p-q,u-v\rangle\ge0$ となる．
これを上の二つの劣勾配に適用すると
\[
 0\le\langle(x-u)-(y-v),u-v\rangle
 =\langle x-y,u-v\rangle-\norm{u-v}^2.
\]
Cauchy--Schwarzの不等式から
$\norm{u-v}\le\norm{x-y}$ が従う．
\end{proof}

\begin{example}[歩長が大きすぎる場合]
$f(x)=Lx^2/2$（$L>0$）とする．
歩長 $\alpha>0$ の勾配法
$x_{k+1}=x_k-\alpha f'(x_k)$ は
\[
 x_k=(1-\alpha L)^kx_0
\]
となる．従って
$\alpha>2/L$ で $x_0\ne0$ なら発散する．
\end{example}
\begin{proof}
$f'(x)=Lx$ なので
$x_{k+1}=(1-\alpha L)x_k$ である．
帰納法から表示式を得る．
$\alpha>2/L$ なら
$|1-\alpha L|>1$ だから
$|x_k|\to\infty$ となる．
\end{proof}

\begin{remark}
勾配法・Newton法・射影勾配法は，
目的関数の平滑性，強凸性，制約集合の幾何といった
仮定に依存して収束の速さが異なる．
収束率の主張には，これらの定量的仮定を
明示することが不可欠である．
\end{remark}

\section{到達課題の完全解答}
\noindent\textbf{勾配法の刻み幅の限界．}
\begin{proposition}
$f(x)=x^2/2$ に一定刻み幅 $\alpha>0$ の
勾配法を適用すると
$x_{n+1}=(1-\alpha)x_n$ となる．
任意の初期値 $x_0$ に対し $x_n\to0$ となる
必要十分条件は $0<\alpha<2$ である．
\end{proposition}
\begin{proof}
$f'(x)=x$ なので
$x_{n+1}=x_n-\alpha f'(x_n)=(1-\alpha)x_n$．
帰納法で $x_n=(1-\alpha)^nx_0$ となる．
全ての $x_0$ に対し収束するためには
$|1-\alpha|<1$ が必要十分であり，
これは $0<\alpha<2$ と同値である．
$\alpha=2$ では $x_n=(-1)^nx_0$ となり，
非零初期値で収束しない．
\end{proof}
\noindent\textbf{到達確認．}
$\alpha=1/2$，$x_0=8$ のとき $x_3$ を求めよ．
\par\noindent\textbf{解答．}
$x_3=(1-1/2)^3\cdot8=1$ である．
\end{document}